% Gap lemmas: scaled gap law, three-distance theorem in the N-point convention, corollaries.
% Labels use the prefix gap:. Classical inputs: continued fraction facts (Khinchin).
\subsection{The scaled half-gap law}\label{gap:sec:mu}

For $N\ge2$ let $g_{N,0},\dots,g_{N,N-1}>0$ be the circular gaps (in radians, summing to $2\pi$) of
$\{\theta_j\}_{j<N}$. The \emph{scaled half-gap law} is the probability measure on $(0,\infty)$
\[
\mu_N=\frac1N\sum_{j=0}^{N-1}\delta_{\,N g_{N,j}/2},
\]
with atoms $N g_{N,j}/2$ and weights $1/N$ (atoms at equal gaps add).

\begin{lemma}\label{gap:lem:mu}
For all $s\ge0$,
\[
\int\ell\,d\mu_N(\ell)=\pi
\qquad\text{and}\qquad
\pi\,m_N(s)=\int(\ell-s)_+\,d\mu_N(\ell).
\]
\end{lemma}
\begin{proof}
By the definition of $\mu_N$ and by \eqref{eq:mN},
\begin{align*}
\int\ell\,d\mu_N&=\frac1N\sum_j\frac N2g_{N,j}=\pi,\\
\pi m_N(s)&=\frac12\sum_j\Bigl(g_{N,j}-\frac{2s}N\Bigr)_+
=\frac1N\sum_j\Bigl(\frac N2g_{N,j}-s\Bigr)_+ .\qedhere
\end{align*}
\end{proof}

\subsection{The three-distance theorem in the $N$-point convention}\label{gap:sec:three}

Let $\beta\in(0,\tfrac12)$ be irrational, with continued fraction $\beta=[0;a_1,a_2,\dots]$.
The convergent denominators and numerators are defined by
\[
q_{-1}=0,\quad q_0=1,\quad q_{k+1}=a_{k+1}q_k+q_{k-1};\qquad
p_{-1}=1,\quad p_0=0,\quad p_{k+1}=a_{k+1}p_k+p_{k-1}.
\]
The signed and absolute errors of the convergents are
\[
\delta_k=q_k\beta-p_k,\qquad \eta_k=|\delta_k|,
\]
so that $\eta_{-1}=1$ and $\eta_0=\beta$. We use the classical facts
(Khinchin, \emph{Continued Fractions}, \S\S1--2):
\begin{align}
q_{k-1}p_k-q_kp_{k-1}&=(-1)^{k-1},\label{gap:eq:det}\\
\delta_k&=(-1)^k\eta_k,\label{gap:eq:sign}\\
\eta_{k+1}&=\eta_{k-1}-a_{k+1}\eta_k>0,\label{gap:eq:rec}\\
q_k\eta_{k-1}+q_{k-1}\eta_k&=1\qquad(k\ge0).\label{gap:eq:F2}
\end{align}
The last identity follows from the first two: by \eqref{gap:eq:det},
\[
q_k\delta_{k-1}-q_{k-1}\delta_k=q_{k-1}p_k-q_kp_{k-1}=(-1)^{k-1},
\]
and $\delta_{k-1}$, $\delta_k$ have opposite signs by \eqref{gap:eq:sign}.
Since $a_1\ge2$,
\[
q_{k-1}<q_k\quad(k\ge1),\qquad \eta_k\le\eta_0<\tfrac12\quad(k\ge0).
\]
Convention: gaps are those of the $N$ points $\{j\beta\bmod1\}_{0\le j<N}$ on $\mathbb R/\mathbb Z$;
the gaps of $\{\theta_j\}$ are $2\pi$ times these.

\begin{theorem}[Three distances, $N$ points]\label{gap:thm:three}
Let $N\ge2$. There is a unique triple $(k,r,t)$ with
\begin{equation}\label{gap:eq:decomp}
N=rq_k+q_{k-1}+t,\qquad k\ge0,\quad 1\le r\le a_{k+1},\quad 0\le t<q_k;
\end{equation}
namely, $k$ is the index with
\[
q_k+q_{k-1}\le N<q_{k+1}+q_k,
\]
and $(r,t)$ is the quotient and remainder of $N-q_{k-1}$ divided by $q_k$.
The gaps of $\{j\beta\}_{j<N}$ are
\begin{align*}
A&=\eta_k &&(N-q_k\text{ times}),\\
B&=\eta_{k+1}+(a_{k+1}-r)\eta_k=\eta_{k-1}-r\eta_k &&(t\text{ times}),\\
C&=A+B=\eta_{k-1}-(r-1)\eta_k &&(q_k-t\text{ times}),
\end{align*}
where for $t=0$ class $B$ is empty. Moreover
\begin{equation}\label{gap:eq:AB}
(N-t)A+q_kB=1.
\end{equation}
\end{theorem}

\emph{Idea of the proof.} Steps 1 and 2 are the classical argument: the closest forward and backward returns $\mathrm d_\pm\le N-1$ determine all gaps, and these returns are $q_k$ and $\rho q_k+q_{k-1}$ for the right $\rho$. Step 3 matches the two descriptions and treats $t=0$ separately.

\begin{proof}
\emph{Preliminaries: uniqueness and counts.} The intervals
\[
I_k=[\,q_k+q_{k-1},\,q_{k+1}+q_k),\qquad k\ge0,
\]
partition $[1,\infty)$, because $q_0+q_{-1}=1$ and the right end of $I_k$ is the left end of $I_{k+1}$.
For $N\in I_k$,
\[
N-q_{k-1}\in[\,q_k,\,(a_{k+1}+1)q_k),
\]
which gives $(r,t)$ uniquely. The counts are nonnegative and sum to $N$, since
\[
(N-q_k)+t+(q_k-t)=N.
\]
For \eqref{gap:eq:AB}, write $N-t=rq_k+q_{k-1}$; then by \eqref{gap:eq:F2},
\[
(N-t)\eta_k+q_k(\eta_{k-1}-r\eta_k)=q_{k-1}\eta_k+q_k\eta_{k-1}=1.
\]
The two forms of $B$ agree by \eqref{gap:eq:rec}.

\emph{Preliminaries: the case $k=0$.} Then $N=r\le a_1$, $t=0$ and $q_0=1$. The points are ordered,
\[
(N-1)\beta<a_1\beta<1,\qquad\text{so}\qquad 0<\beta<\dots<(N-1)\beta<1,
\]
with gaps
\begin{align*}
\beta&=\eta_0 &&(N-1\text{ times}),\\
1-(N-1)\beta&=C &&(\text{once}).
\end{align*}

\emph{Step 1: general structure.} Let $N\ge2$ and let $\mathrm d_+$ (resp.\ $\mathrm d_-$) be the $d\in[1,N-1]$ with
$\{d\beta\}$ minimal (resp.\ maximal); these are unique since the $\{d\beta\}$ are distinct. Put
\[
\Delta_+=\{\mathrm d_+\beta\},\qquad \Delta_-=1-\{\mathrm d_-\beta\}.
\]
Then:
\begin{enumerate}
\item[(a)] $\mathrm d_++\mathrm d_-\ge N$. Otherwise
\[
d=\mathrm d_++\mathrm d_-\le N-1
\qquad\text{and}\qquad
\{d\beta\}\equiv\Delta_+-\Delta_-\ne0,
\]
so
\[
\{d\beta\}=
\begin{cases}
\Delta_+-\Delta_-\in(0,\Delta_+) & \text{if }\Delta_+>\Delta_-,\\
1-\Delta_-+\Delta_+>\{\mathrm d_-\beta\} & \text{if }\Delta_+<\Delta_-.
\end{cases}
\]
Both contradict the choice of $\mathrm d_\pm$.
\item[(b)] The gaps are
\begin{align*}
&\Delta_+ &&(N-\mathrm d_+\text{ times}),\\
&\Delta_- &&(N-\mathrm d_-\text{ times}),\\
&\Delta_++\Delta_- &&(\mathrm d_++\mathrm d_--N\text{ times}).
\end{align*}
To prove this, let $\operatorname{nxt}(j)$ be the successor of the point $j\beta$ ($0\le j<N$) in the
counterclockwise order; its index is $j+d$, where $d\ne0$ minimizes the forward distance $\{d\beta\}$ over
$-j\le d\le N-1-j$. Positive $d$ give $\{d\beta\}\ge\Delta_+$; negative $d=-e$ give
\[
\{-e\beta\}=1-\{e\beta\}\ge\Delta_-.
\]
We determine $\operatorname{nxt}(j)$ on three ranges of $j$.
 \begin{itemize}
 \item \emph{Range $j<N-\mathrm d_+$.} Here $\operatorname{nxt}(j)=j+\mathrm d_+$, with gap $\Delta_+$.
 Positive $d$ are fine. For negative $d=-e$ with $e\le j$, suppose $1-\{e\beta\}<\Delta_+$.
 Since $e\le j<N-\mathrm d_+$,
 \[
 e+\mathrm d_+\le N-1
 \qquad\text{and}\qquad
 \{(e+\mathrm d_+)\beta\}=\{e\beta\}+\Delta_+-1\in(0,\Delta_+),
 \]
 which contradicts the choice of $\mathrm d_+$.
 \item \emph{Range $j\ge\mathrm d_-$.} By (a), $j\ge N-\mathrm d_+$, so this range is disjoint from the previous one.
 Here $\operatorname{nxt}(j)=j-\mathrm d_-$, with gap $\Delta_-$. Negative $d$ are fine.
 For positive $d$, suppose $\{d\beta\}<\Delta_-$. Since $d\le N-1-j$ and $j\ge\mathrm d_-$,
 \[
 d+\mathrm d_-\le N-1
 \qquad\text{and}\qquad
 \{(d+\mathrm d_-)\beta\}=\{d\beta\}+1-\Delta_-\in(\{\mathrm d_-\beta\},1),
 \]
 which contradicts the choice of $\mathrm d_-$.
 \item \emph{Range $N-\mathrm d_+\le j<\mathrm d_-$.} Here $\operatorname{nxt}(j)=j+\mathrm d_+-\mathrm d_-$, with gap $\Delta_++\Delta_-$.
 The index lies in $[0,N-1]$, because
 \[
 j+\mathrm d_+-\mathrm d_-\ge N-\mathrm d_-\ge1
 \qquad\text{and}\qquad
 j+\mathrm d_+-\mathrm d_-<\mathrm d_+\le N-1.
 \]
 Positive available $d$ satisfy $d\le\mathrm d_+-1$. For such $d$, suppose $\Delta_+<\{d\beta\}<\Delta_++\Delta_-$.
 Then $e=\mathrm d_+-d\in[1,\mathrm d_+-1]$ and
 \[
 \{e\beta\}=1-(\{d\beta\}-\Delta_+)>1-\Delta_-,
 \]
 which contradicts the choice of $\mathrm d_-$.
 For negative $d=-e$ with $e\le j<\mathrm d_-$, suppose $\Delta_-<1-\{e\beta\}<\Delta_++\Delta_-$.
 Then $d'=\mathrm d_--e\in[1,\mathrm d_--1]$ and
 \[
 \{d'\beta\}=1-\Delta_--\{e\beta\}\in(0,\Delta_+),
 \]
 which contradicts the choice of $\mathrm d_+$.

 Finally, the candidate index itself has forward distance $\Delta_++\Delta_-$: indeed
 \[
 (\mathrm d_+-\mathrm d_-)\beta\equiv\Delta_++\Delta_-\pmod1,
 \]
 and $\Delta_++\Delta_-<1$. For the latter, if $\Delta_++\Delta_-\ge1$, the same argument with $d=\mathrm d_+-\mathrm d_-$ or $e=\mathrm d_--\mathrm d_+$ in $[1,N-1]$ contradicts minimality or maximality; equality would give $(\mathrm d_+-\mathrm d_-)\beta\in\mathbb Z$, hence $\mathrm d_+=\mathrm d_-$ and $\Delta_++\Delta_-=1$, which forces $N=2$, where the third range is empty.
 \end{itemize}
Distinct points give strict inequalities, so these are the successors. The three ranges of $j$ have sizes
$N-\mathrm d_+$, $N-\mathrm d_-$ and $\mathrm d_++\mathrm d_--N$, which gives the multiplicities.
\end{enumerate}

\emph{Step 2: identifying $\mathrm d_\pm$.} Let $k\ge1$ and $0\le\rho\le a_{k+1}$, and put
\[
\mathrm d_\rho=\rho q_k+q_{k-1}\ge1,\qquad
\nu=\eta_{k-1}-\rho\eta_k>0,\qquad
\sigma=(-1)^k;
\]
by \eqref{gap:eq:rec}, $\nu\ge\eta_{k+1}$. Then
\begin{equation}\label{gap:eq:step2}
q_k\beta\equiv\sigma\eta_k,\qquad \mathrm d_\rho\beta\equiv-\sigma\nu\pmod1,\qquad
q_k\nu+\mathrm d_\rho\eta_k=1,
\end{equation}
the last by \eqref{gap:eq:F2}.

\emph{Claim.} If $N\le\mathrm d_\rho+q_k$, there are no $d\in[1,N-1]$ and $p\in\mathbb Z$ with
\[
e=\sigma(d\beta-p)\in(-\nu,\eta_k)\setminus\{0\}.
\]
To see this, consider the vectors
\[
b_1=(q_k,\sigma\eta_k),\qquad b_2=(\mathrm d_\rho,-\sigma\nu).
\]
They have determinant $-\sigma$ and span the lattice $\mathcal L=\{(d,d\beta-p):d,p\in\mathbb Z\}$;
unimodularity follows from \eqref{gap:eq:det}, since $(\mathrm d_\rho,\cdot)=\rho(q_k,\cdot)+(q_{k-1},\cdot)$.
Write $(d,d\beta-p)=mb_1+nb_2$ with $m,n\in\mathbb Z$. Cramer's rule gives
\[
m=d\nu+\mathrm d_\rho e,\qquad n=\eta_kd-q_ke,\qquad e=\sigma(d\beta-p).
\]
Suppose first $e\in(0,\eta_k)$. Then $m\ge1$, and
\[
d=mq_k+n\mathrm d_\rho\le N-1<q_k+\mathrm d_\rho
\]
forces $n\le0$, so
\[
e=m\eta_k-n\nu\ge\eta_k,
\]
a contradiction. Suppose next $e\in(-\nu,0)$. Then $n\ge1$, the same bound on $d$ forces $m\le0$, and
\[
e=m\eta_k-n\nu\le-\nu,
\]
a contradiction. This proves the claim.

Consequently, for $d\in[1,N-1]$,
\[
\{d\beta\}\notin
\begin{cases}
(0,\eta_k)\cup(1-\nu,1) & \text{if }\sigma=1,\\
(0,\nu)\cup(1-\eta_k,1) & \text{if }\sigma=-1.
\end{cases}
\]
If also $q_k\le N-1$ and $\mathrm d_\rho\le N-1$, the values attained there are extremal:
\[
\bigl\{\{q_k\beta\},\{\mathrm d_\rho\beta\}\bigr\}=
\begin{cases}
\{\eta_k,\,1-\nu\} & \text{if }\sigma=1,\\
\{1-\eta_k,\,\nu\} & \text{if }\sigma=-1.
\end{cases}
\]
Hence
\[
\{\mathrm d_+,\mathrm d_-\}=\{q_k,\mathrm d_\rho\},\qquad \{\Delta_+,\Delta_-\}=\{\eta_k,\nu\},
\]
and by Step 1(b) the gaps are
\begin{align*}
&\eta_k &&(N-q_k\text{ times}),\\
&\nu &&(N-\mathrm d_\rho\text{ times}),\\
&\eta_k+\nu &&(q_k+\mathrm d_\rho-N\text{ times}).
\end{align*}

\emph{Step 3: apply with $k\ge1$.} Let $(k,r,t)$ be as in \eqref{gap:eq:decomp} with $k\ge1$.
If $t\ge1$, take $\rho=r$. Then
\begin{align*}
\mathrm d_r&=N-t\le N-1, & N-\mathrm d_r&=t,\\
q_k&\le N-1\quad(\text{as }N\ge q_k+q_{k-1}+1), & N&=\mathrm d_r+t\le\mathrm d_r+q_k,
\end{align*}
so Step 2 applies with $\nu=B$, and the counts are $N-q_k$, $t$, $q_k-t$ for $A$, $B$, $C$.
If $t=0$, take $\rho=r-1$. Then
\begin{align*}
\mathrm d_\rho&=N-q_k\ge q_{k-1}\ge1, & N&=\mathrm d_\rho+q_k,\\
q_k&\le N-1, & \nu&=\eta_{k-1}-(r-1)\eta_k=C,
\end{align*}
and the counts are $N-q_k$ for $A$, $q_k$ for $C$, and $0$ for $C+A$. This is the claim with $t=0$.
\end{proof}

\begin{remark}[Relation to Leong]
Leong \cite[Thm.~3]{Leong17} treats the $N'+1$ points $0,\beta,\dots,N'\beta$, i.e.\ $N=N'+1$.
Under that shift the ranges
\[
q_k+q_{k-1}\le N'<q_{k+1}+q_k
\qquad\text{become}\qquad
q_k+q_{k-1}<N\le q_{k+1}+q_k,
\]
so the boundary values $N=q_k+q_{k-1}$ (that is, $r=1$, $t=0$) fall at the end of the previous level's range.
A direct citation therefore needs the level drop $k\to k-1$ (with $r'=a_k$) for $r=1$ and for $N=q_k$.
Theorem~\ref{gap:thm:three} avoids this: every $N\ge2$ is covered exactly once with $t=0$ allowed,
and its proof is self-contained apart from \eqref{gap:eq:det}--\eqref{gap:eq:F2}.
\end{remark}

\subsection{Corollaries}\label{gap:sec:cor}

For $k\ge0$ write
\begin{equation}\label{gap:eq:xi}
\xi_k=\frac{\eta_k}{\eta_{k-1}},
\qquad\text{so that}\qquad
\xi_k=\frac1{a_{k+1}+\xi_{k+1}}
\end{equation}
by \eqref{gap:eq:rec}. Moreover $\xi_k\in(0,1)$, because
\[
\eta_{k-1}\ge a_{k+1}\eta_k+\eta_{k+1}>\eta_k.
\]

\begin{corollary}[Standard inequalities]\label{gap:cor:ineq}
For $k\ge0$:
\begin{enumerate}
\item[(i)] $q_k\eta_{k-1}+q_{k-1}\eta_k=1$;
\item[(ii)] the two-sided bound
\[
\frac1{a_{k+1}+2}<q_k\eta_k<\frac1{a_{k+1}}.
\]
\end{enumerate}
\end{corollary}
\begin{proof}
Part (i) is \eqref{gap:eq:F2}. For (ii), apply \eqref{gap:eq:F2} at $k+1$ and divide by $q_k\eta_k$:
\[
\frac1{q_k\eta_k}=\frac{q_{k+1}}{q_k}+\frac{\eta_{k+1}}{\eta_k}.
\]
Here
\[
a_{k+1}\le\frac{q_{k+1}}{q_k}=a_{k+1}+\frac{q_{k-1}}{q_k}<a_{k+1}+1
\qquad\text{and}\qquad
0<\frac{\eta_{k+1}}{\eta_k}<1.\qedhere
\]
\end{proof}

\begin{corollary}[Two-distance times]\label{gap:cor:two}
Two cases of Theorem~\ref{gap:thm:three} have only two gap lengths.
\begin{enumerate}
\item[(i)] Let $N=rq_k+q_{k-1}$ with $k\ge0$, $1\le r\le a_{k+1}$ and $N\ge2$. Then the gaps are
\begin{align*}
&\eta_k &&(N-q_k\text{ times}),\\
&\eta_{k-1}-(r-1)\eta_k &&(q_k\text{ times}).
\end{align*}
\item[(ii)] Let $N=q_k\ge2$ with $k\ge1$. Then the gaps are
\begin{align*}
&\eta_{k-1} &&(q_k-q_{k-1}\text{ times}),\\
&\eta_{k-1}+\eta_k &&(q_{k-1}\text{ times}).
\end{align*}
\end{enumerate}
\end{corollary}
\begin{proof}
(i) is Theorem~\ref{gap:thm:three} with $t=0$; the decomposition of $N$ is $(k,r,0)$.

(ii) Since
\[
N=q_k=a_kq_{k-1}+q_{k-2},
\]
the decomposition of $N$ is $(k-1,\,r=a_k,\,t=0)$. Part (i) at level $k-1$ gives the gaps $\eta_{k-1}$ ($q_k-q_{k-1}$ times) and
\[
\eta_{k-2}-(a_k-1)\eta_{k-1}=\eta_k+\eta_{k-1}\qquad(q_{k-1}\text{ times}),
\]
using \eqref{gap:eq:rec} in the form $\eta_k=\eta_{k-2}-a_k\eta_{k-1}$.
For $k=1$ use the case $k=0$ of the theorem ($q_1=a_1$).
\end{proof}
These are the cases used in Theorem~\ref{thm:B}(2). There, by Corollary~\ref{gap:cor:ineq},
\begin{equation}\label{gap:eq:twoscaled}
\begin{aligned}
q_k\eta_{k-1}&=1-q_{k-1}\eta_k\in\Bigl(1-\frac1{a_{k+1}},\,1\Bigr),\\
q_k(\eta_{k-1}+\eta_k)&=1+(q_k-q_{k-1})\eta_k\in\Bigl(1,\,1+\frac1{a_{k+1}}\Bigr).
\end{aligned}
\end{equation}

\begin{corollary}[Bounded type]\label{gap:cor:bdd}
Let $a_j\le A_0$ for all $j$, and put
\[
\varepsilon=\frac1{(A_0+1)(A_0+2)}.
\]
Let $k\ge1$ and let $N=rq_k+q_{k-1}+t$ be as in Theorem~\ref{gap:thm:three}. Then:
\begin{enumerate}
\item[(i)] the ratio $\eta_{k+1}/\eta_k$ satisfies
\[
\frac1{A_0+1}<\frac{\eta_{k+1}}{\eta_k}<\frac{A_0+1}{A_0+2};
\]
\item[(ii)] the class frequencies satisfy
\[
\frac{N-q_k}N\ge\varepsilon\qquad\text{and}\qquad\frac{q_k}N\ge\frac1{A_0+2};
\]
\item[(iii)] the gap ratios satisfy
\begin{align*}
C/A&\ge1+1/(A_0+1) &&\text{always},\\
B/A&\ge1+1/(A_0+1) &&\text{if }r<a_{k+1},\\
A/B&\ge1+1/(A_0+1) &&\text{if }r=a_{k+1}.
\end{align*}
\end{enumerate}
\end{corollary}
\begin{proof}
(i) By \eqref{gap:eq:xi},
\[
\xi_{k+1}=\frac1{a_{k+2}+\xi_{k+2}}>\frac1{A_0+1}
\qquad\text{and}\qquad
\xi_{k+1}<\frac1{1+\xi_{k+2}}<\frac1{1+1/(A_0+1)}.
\]

(ii) We have
\[
N-q_k=(r-1)q_k+q_{k-1}+t\ge q_{k-1}\ge\frac{q_k}{A_0+1}
\qquad\text{and}\qquad
N<q_{k+1}+q_k\le(A_0+2)q_k.
\]
Hence
\[
\frac{N-q_k}N\ge\frac{q_{k-1}}{(A_0+2)q_k}\ge\varepsilon
\qquad\text{and}\qquad
\frac{q_k}N>\frac1{A_0+2}.
\]

(iii) First, $C/A=1+B/A$. If $r<a_{k+1}$, then $B\ge\eta_{k+1}+\eta_k$, so by (i)
\[
\frac BA\ge1+\xi_{k+1}>1+\frac1{A_0+1}
\qquad\text{and}\qquad
\frac CA>2.
\]
If $r=a_{k+1}$, then $B=\eta_{k+1}$, so by (i)
\[
\frac AB=\frac1{\xi_{k+1}}>1+\frac1{A_0+1}
\qquad\text{and}\qquad
\frac CA=1+\xi_{k+1}>1+\frac1{A_0+1}.\qedhere
\]
\end{proof}
This is exactly what the proof of Theorem~\ref{thm:B}(3) needs, now
with the unified decomposition valid for all $N\ge2$; for $k=0$ the frequency of class $A$ may be $0$, which affects only the finitely many $N\le a_1$.

\begin{corollary}[Golden law]\label{gap:cor:golden}
Let
\[
\lambda=\frac{\sqrt5-1}2,\qquad F_1=F_2=1,\qquad \beta=\lambda^2=[0;2,1,1,\dots].
\]
Let $m\ge3$ and $N=F_m+t$ with $0\le t<F_{m-1}$. Then the gaps of $\{j\beta\}_{j<N}$ are
\begin{align*}
&\lambda^{m-2} &&(F_{m-1}-t\text{ times}),\\
&\lambda^{m-1} &&(F_{m-2}+t\text{ times}),\\
&\lambda^{m} &&(t\text{ times}).
\end{align*}
\end{corollary}
\begin{proof}
Here $a_1=2$ and $a_j=1$ for $j\ge2$. For $j\ge0$,
\begin{equation}\label{gap:eq:goldid}
q_j=F_{j+2},\qquad \eta_j=\lambda^{j+2}.
\end{equation}
Indeed $q_0=1$, $q_1=2$, and
\[
\eta_0=\lambda^2,\qquad \eta_1=1-2\lambda^2=\lambda^3,
\]
while the recursion $\eta_{j+1}=\eta_{j-1}-\eta_j$ matches $\lambda^{j+3}=\lambda^{j+1}-\lambda^{j+2}$.

For $m=3$, $N=2$ is the case $k=0$, $r=2$ of the theorem, with gaps
\[
\lambda^2,\qquad 1-\lambda^2=\lambda,
\]
matching the multiplicities $(F_2-0,\,F_1+0,\,0)=(1,1,0)$.

For $m\ge4$ put $K=m-3\ge1$. Then
\[
q_K+q_{K-1}=F_{m-1}+F_{m-2}=F_m\le N<F_{m+1}=q_{K+1}+q_K,
\]
so $N$ lies at level $K$. Since $a_{K+1}=1$, we get $r=1$ and $t=N-F_m<q_K=F_{m-1}$.
The theorem gives
\begin{align*}
A&=\eta_K=\lambda^{m-1} &&(N-q_K=F_{m-2}+t\text{ times}),\\
B&=\eta_{K+1}=\lambda^{m} &&(t\text{ times}),\\
C&=\eta_{K-1}=\lambda^{m-2} &&(F_{m-1}-t\text{ times}).\qedhere
\end{align*}
\end{proof}
Equivalently, one may apply the theorem to $\lambda=[0;1,1,\dots]$ after the reflection $x\mapsto-x$, which maps $\{j\lambda\}$ to $\{j\beta\}$ since $\beta\equiv-\lambda$; indices then shift by one. The direct formulation above avoids $a_1=1$, which is outside our range $\beta<\tfrac12$.

\begin{corollary}[Noble law]\label{gap:cor:noble}
Let $a_j=1$ for all $j\ge k_0$.
\begin{enumerate}
\item[(i)] For $k>k_0$ and $N=q_k+t$ with $0\le t<q_{k-1}$, the gaps are
\begin{align*}
&\eta_{k-1} &&(N-q_{k-1}\text{ times}),\\
&\eta_k &&(t\text{ times}),\\
&\eta_{k-2} &&(q_{k-1}-t\text{ times}).
\end{align*}
\item[(ii)] For $k\ge k_0$, $\eta_k/\eta_{k-1}=\lambda$. Moreover, as $k\to\infty$,
\[
\frac{q_{k-1}}{q_k}\to\lambda,\qquad
q_k\eta_{k-1}\to\frac1{1+\lambda^2}=\frac1{\sqrt5\,\lambda},\qquad
q_k\eta_{k-2}\to\frac1{\sqrt5\,\lambda^2}.
\]
\end{enumerate}
\end{corollary}
\begin{proof}
(i) Put $K=k-1\ge k_0\ge1$. Since $a_k=1$, $q_k=q_K+q_{K-1}$, hence
\[
N=q_K+q_{K-1}+t,\qquad r=1=a_{K+1},\qquad 0\le t<q_K,
\]
so the decomposition of $N$ is $(K,1,t)$. Theorem~\ref{gap:thm:three} gives
\begin{align*}
A&=\eta_K=\eta_{k-1} &&(N-q_{k-1}\text{ times}),\\
B&=\eta_{K+1}=\eta_k &&(t\text{ times}),\\
C&=\eta_{K-1}=\eta_{k-2} &&(q_{k-1}-t\text{ times}).
\end{align*}

(ii) For $j\ge k_0$, \eqref{gap:eq:xi} reads $\xi_j=f(\xi_{j+1})$, where
\[
f(x)=\frac1{1+x},\qquad |f(x)-\lambda|\le\lambda\,|x-\lambda|\quad\text{on }[0,\infty).
\]
Iterating, $|\xi_j-\lambda|\le\lambda^n$ for every $n$, so $\xi_j=\lambda$.
The same contraction applied to
\[
\frac{q_{j-1}}{q_j}=f\Bigl(\frac{q_{j-2}}{q_{j-1}}\Bigr)
\]
gives $q_{k-1}/q_k\to\lambda$. Finally, by \eqref{gap:eq:F2} and \eqref{gap:eq:xi},
\[
q_k\eta_{k-1}=\Bigl(1+\frac{q_{k-1}}{q_k}\,\xi_k\Bigr)^{-1}\to\frac1{1+\lambda^2}=\frac1{\sqrt5\,\lambda},
\qquad
\eta_{k-2}=\frac{\eta_{k-1}}{\lambda},
\]
where $1+\lambda^2=\sqrt5\,\lambda$.
\end{proof}
Consequently, writing
\[
x_N=\frac t{q_k}\in\Bigl[0,\frac{q_{k-1}}{q_k}\Bigr)
\]
and using $\eta_{k-1}=\lambda\eta_{k-2}$, $\eta_k=\lambda^2\eta_{k-2}$, the measure $\mu_N$ has
\begin{align*}
\text{atoms}&\quad \pi N\eta_{k-2}\cdot\{1,\lambda,\lambda^2\},\\
\text{weights}&\quad \frac{q_{k-1}-t}N,\ \frac{N-q_{k-1}}N,\ \frac tN,
\end{align*}
respectively, and
\[
N\eta_{k-2}=(1+x_N)\,q_k\eta_{k-2}\to\frac{1+x}{\sqrt5\,\lambda^2}.
\]
This is the family $\mu_x$ of Section~\ref{sec:golden}.
